\begin{textsample}{POS Dim 2 – pe – Score 18.00 – pe/pe\_em\_21.txt}  \label{ex:f2_pos_275}
2.8 – \textbf{ITINERÁRIOS} FORMATIVOS Os \textbf{Itinerários} Formativos, segundo as DCNEM, são : [... ] cada conjunto de unidades curriculares \textbf{ofertadas} pelas \textbf{instituições} e redes de ensino que possibilitam ao estudante aprofundar seus conhecimentos e se preparar para o \textbf{prosseguimento} de estudos ou para o mundo do trabalho de forma a contribuir para a construção de soluções de problemas específicos da sociedade ( BRASIL, 2018 ).
Segundo o mesmo documento, são cinco os \textbf{itinerários} formativos : Linguagens e suas tecnologias, Matemática e suas Tecnologias, Ciências da Natureza e suas Tecnologias, Ciências Humanas e Sociais Aplicadas e Formação \textbf{Técnica} e Profissional.
Os \textbf{Itinerários} Formativos são a parte do \textbf{currículo} no qual os estudantes deverão optar por uma trilha de aprofundamento, em uma das áreas de conhecimento, ou em áreas de conhecimento integradas, conforme tenham afinidade e/ou interesse, ou por uma trilha da formação \textbf{técnica} e profissional.
Os \textbf{Itinerários} pretendem assim promover o aprofundamento e a ampliação das aprendizagens dos estudantes em relação à FGB.
Em consonância com as DCNEM, o Ministério da Educação ( MEC ) estabeleceu os Referenciais para Elaboração dos \textbf{Itinerários} Formativos, através da \textbf{Portaria} nº 1.432/2018, que estabelece que os \textbf{itinerários} devem ser organizados em torno de um ou mais dos eixos estruturantes a seguir : Investigação Científica, Processos Criativos, Mediação e Intervenção Sociocultural e Empreendedorismo. [... ] I - investigação científica : supõe o aprofundamento de conceitos fundantes das ciências para a interpretação de ideias, fenômenos e processos para serem utilizados em procedimentos de investigação voltados ao enfrentamento de situações cotidianas e demandas locais e coletivas, e a proposição de intervenções que considerem o desenvolvimento local e a melhoria da qualidade de vida da comunidade ; II - processos criativos : supõe o uso e o aprofundamento do conhecimento científico na construção e criação de experimentos, modelos, protótipos para a criação de processos ou produtos que \textbf{atendam} a demandas pela resolução de problemas identificados na sociedade ; III - mediação e intervenção sociocultural : supõe a mobilização de conhecimentos de uma ou mais áreas para mediar conflitos, promover entendimento e implementar soluções para questões e problemas identificados na comunidade ; IV - empreendedorismo : supõe a mobilização de conhecimentos de diferentes áreas para a formação de organizações com variadas missões voltadas ao desenvolvimento de produtos ou prestação de serviços inovadores com o uso das tecnologias. ( Resolução CNE/CEB nº 3/2018, Art. 12, § 2º ) Nesta perspectiva, estes eixos permearão todos os \textbf{Itinerários} Formativos propostos no \textbf{Currículo} de Pernambuco para o Ensino Médio, por entender que as habilidades que estão relacionadas a eles contribuem, significativamente, para a formação dos estudantes.
Os \textbf{Itinerários} Formativos, dentro desse novo \textbf{currículo}, têm ainda o papel de contribuir para que os estudantes desenvolvam a autonomia de forma que tenham mais elementos para a realização de seus projetos de vida.
Os eixos estruturantes e suas habilidades ajudam a promover uma visão de mundo mais ampla dos estudantes e, assim, contribuem para que eles possam tomar decisões, almejando uma formação para além da escola, pessoal e profissional.
A \textbf{carga} \textbf{horária} \textbf{destinada} aos \textbf{Itinerários} Formativos está distribuída em tipos diferentes de unidades curriculares que são : obrigatórias, optativas, \textbf{eletivas} e projeto de vida.
As Unidades Obrigatórias são aquelas que têm o papel de aprofundamento dentro da área/temática de escolha dos estudantes para o percurso do Ensino Médio e, como o nome indica, são obrigatórias para todos os estudantes que resolveram \textbf{cursar} uma determinada trilha.
Nessas unidades serão abordados conceitos e conhecimentos que possibilitam aos estudantes refletirem mais atentamente sobre um problema epistemológico, social, ambiental, cultural, de uma ou mais áreas de conhecimento em que tenham interesse.
As Unidades Optativas são aquelas que irão compor, junto com as obrigatórias, o percurso formativo relacionado à área ou áreas de conhecimento ; deverão, então, ser escolhidas pela escola a partir de um \textbf{catálogo} de opções disponibilizado pela Secretaria de Educação e Esportes.
As Unidades Curriculares \textbf{Eletivas}, por sua vez, são aquelas que visam ampliar o universo de conhecimentos dos estudantes, em seus interesses mais diversos.
Necessariamente, não precisam estar diretamente relacionadas à área de conhecimento escolhida pelo estudante.
Essas \textbf{eletivas} serão propostas pela escola, em articulação com o interesse do educando e a formação dos professores, com acompanhamento da SEE.
Dessa forma, podem ser explorados conhecimentos sobre diversos temas, desde que \textbf{atendam} aos critérios acima descritos — interesse do educando e a formação dos professores — e corroborem para a formação dos estudantes.
É importante salientar que dá -se ao estudante autonomia, nesse processo, para escolher qual \textbf{eletiva} \textbf{cursará}.
A Unidade Curricular de Projeto de Vida busca despertar nos estudantes uma reflexão sobre o seu futuro pessoal e profissional, incluindo elementos relativos ao autoconhecimento, ao conhecimento do outro e ao papel que todos temos que desempenhar na sociedade em que vivemos.
Nesse processo, discussões e estudos pretendem incentivar o estudante a refletir sobre sua vida, estimulando -o a fazer escolhas e tomar decisões, haja vista a concretização de suas expectativas.
Diferente das unidades curriculares na Formação Geral Básica, trabalhadas anualmente, as unidades curriculares dos \textbf{Itinerários} Formativos terão a duração de um \textbf{semestre}, possibilitando a vivência de uma gama maior de temas e experiências ao longo do Ensino Médio.
Para cada um dos \textbf{Itinerários} Formativos indicados nas DCNEM, a fim de ampliar a possibilidade de escolha dos estudantes, o \textbf{Currículo} de Pernambuco para o Ensino Médio está \textbf{ofertando} de uma a três Trilhas de aprofundamento específica por área de conhecimento, além das trilhas que integram duas áreas de conhecimento.
Essas trilhas específicas de área não estão ilhadas em si mesmas ; antes, estabelecem diálogo com outras áreas, embora contemplem predominantemente unidades curriculares associadas à própria área.
Como dito anteriormente, também são \textbf{ofertados} \textbf{Itinerários} Integrados que articulam preponderantemente duas áreas de conhecimento, garantindo ao estudante a possibilidade de ampliar seu olhar em relação às temáticas trabalhadas sob diferentes ângulos.
Todas as áreas de conhecimento estão integradas entre si, promovendo um movimento de interdisciplinaridade entre os componentes e, consequentemente, uma formação mais ampla para os estudantes.
A perspectiva de aprofundamento, a partir de \textbf{Itinerários} Formativos, reforça a proposta da interdisciplinaridade, mas também o diálogo entre áreas, pois em todas as trilhas, específicas ou integradas, sempre se encontram unidades curriculares de mais de uma área de conhecimento.
O que difere uma trilha específica de uma integrada, na verdade, é que nestas últimas há maior equilíbrio entre a quantidade de unidades curriculares de diferentes áreas, enquanto nas específicas, os componentes de uma área, a que deu origem à trilha, prevalecem em detrimento de outros.
Considerando que os IF representam grande inovação no \textbf{currículo} do Ensino Médio, o próximo \textbf{capítulo} aborda especialmente os \textbf{Itinerários} Formativos de Áreas de Conhecimento e os \textbf{Itinerários} Formativos da Educação Profissional \textbf{Técnica}

% matched lemmas: atender, capítulo, carga, catálogo, currículo, cursar, destinar, eletivo, horário, instituição, itinerário, ofertar, portaria, prosseguimento, semestre, técnico
\end{textsample}
